% Hodina 29, 1. ročník SPŠ. Rozsah: Rozpis_vyuky_MAT_1_rocnik.tex, řádek 29.

% Plán: 1_rocnik/tematicke_plany/SPS.pdf; společná projektová šablona.

% Učebnice: ucebnice/01_Zakladni_poznatky/Mocniny_a_odmocniny/Mocniny a Odmocniny.tex

% Původ úloh: Vlastní zadání a vlastní formulace výkladu; metodická inspirace zdroji.

% Metodický zdroj: externi_zdroje/realisticky/Mocniny/08_Exponencialni_tvar_cisla.pdf

% URL: http://www.realisticky.cz/ucebnice/01%20Matematika%20S%C5%A0/01%20Z%C3%A1kladn%C3%AD%20poznatky/06%20Mocniny/08%20Exponenci%C3%A1ln%C3%AD%20tvar%20%C4%8D%C3%ADsla.pdf

% Metodický zdroj: externi_zdroje/realisticky/Mocniny/09_Vypocty_s_mocninami.pdf

% URL: http://www.realisticky.cz/ucebnice/01%20Matematika%20S%C5%A0/01%20Z%C3%A1kladn%C3%AD%20poznatky/06%20Mocniny/09%20V%C3%BDpo%C4%8Dty%20s%20mocninami.pdf

\documentclass[11pt,a4paper]{article}

% Spolecny styl pro ucitelske pripravy MAT (Work/MAT)
% Verze 1.1, 2026-10-23
\NeedsTeXFormat{LaTeX2e}
\ProvidesPackage{priprava}[2026/10/23 v1.1 Sablona ucitelskych priprav MAT]

% Zakladni typografie a cestina
\RequirePackage[T1]{fontenc}
\RequirePackage[utf8]{inputenc}
\RequirePackage[czech]{babel}
\RequirePackage{lmodern}
\RequirePackage{microtype}
\RequirePackage{amsmath,amssymb}
\RequirePackage{xcolor}
\RequirePackage{enumitem}
\RequirePackage[most]{tcolorbox}
\RequirePackage{tabularx,booktabs,array}
\RequirePackage{fancyhdr}
\RequirePackage{needspace}
\RequirePackage{geometry}
\RequirePackage[hidelinks]{hyperref}

\geometry{a4paper,top=17mm,bottom=18mm,left=23mm,right=23mm}
\definecolor{notegray}{gray}{0.39}
\setlength{\parindent}{0pt}
\setlength{\parskip}{2.5pt plus 1pt minus 0.5pt}
\setlist[itemize]{topsep=2pt,itemsep=1pt,parsep=0pt,leftmargin=1.45em}
\setlist[enumerate]{topsep=2pt,itemsep=2pt,parsep=0pt,leftmargin=1.65em}
\setcounter{secnumdepth}{0}
\renewcommand{\arraystretch}{1.12}

% Zahlavi: vlevo rocnik | globalni tema, vpravo cislo: nazev hodiny
% Zapati zustava zcela prazdne.
\pagestyle{fancy}
\fancyhf{}
\renewcommand{\headrulewidth}{0.35pt}
\renewcommand{\footrulewidth}{0pt}
\setlength{\headheight}{14pt}
\newcommand{\nastavPripravu}[4]{%
  \fancyhead[L]{\small #1\enspace\textbar\enspace #2}%
  \fancyhead[R]{\small #3: #4}%
}

% Minimalisticke cernobile boxy. Kazda definice zvlast.
\tcbset{commonlesson/.style={enhanced,breakable=false,colback=white,colframe=black,
  colbacktitle=white,coltitle=black,boxrule=0.45pt,
  arc=0.8mm,outer arc=0.8mm,left=3mm,right=3mm,top=1.6mm,bottom=1.6mm,
  toptitle=1mm,bottomtitle=0.7mm,fonttitle=\bfseries,
  attach boxed title to top left={xshift=2.5mm,yshift=-1.15mm},
  boxed title style={colback=white,colframe=white,boxrule=0pt,left=1mm,right=1mm}}}
\newtcolorbox{zadaniBox}[1][]{commonlesson,title={Zadání},#1}
\newtcolorbox{prikladBox}[1][]{commonlesson,title={Příklad},#1}
\newtcolorbox{definiceBox}[1][]{commonlesson,title={Definice},#1}

% Spolecny box pro zadani a jeho reseni, bez mezery mezi castmi.
% Pouziti:
%   \begin{prikladReseniBox}
%     Zadani ulohy.
%     \tcblower
%     \begin{align*} ... \end{align*}
%   \end{prikladReseniBox}
% Spodni cast ma jemny sedy podklad a delici linku; lze v ni sazet
% vice radku matematiky. Obe casti maji jen jeden spolecny ramecek.
\tcbset{solutionlesson/.style={
  commonlesson,
  enhanced jigsaw,
  breakable,
  colbacklower=black!3,
  segmentation style={solid,draw=black!25,line width=0.35pt},
  fontlower=\small\color{notegray},
  before lower={\noindent\textbf{Řešení:}\enspace},
  leftlower=3mm,rightlower=3mm
}}
\newtcolorbox{prikladReseniBox}[1][]{solutionlesson,title={Příklad},#1}
\newtcolorbox{zadaniReseniBox}[1][]{solutionlesson,title={Zadání},#1}

% Pomucky pro autora pripravy: sedive poznamky; \reseni pro starsi pripravy
\newcommand{\poznamka}[1]{\par\begingroup\small\color{notegray}\itshape
  \textnormal{\textbf{Poznámka pro učitele.}}\enspace #1\par\endgroup}
\newcommand{\reseni}[1]{\par\begingroup\small\color{notegray}
  \textbf{Očekávané řešení:}\enspace #1\par\endgroup}
\newcommand{\blok}[1]{\needspace{5\baselineskip}\subsection*{#1}}
\newcommand{\vyrok}[1]{\uv{#1}}


\nastavPripravu{1. ročník}{Mocniny}{29}{Vědecký zápis a odhady}

\begin{document}

\blok{Výklad a souvislosti}

Mocniny desítky umožňují porovnat rozměry, které se liší mnoha nulami. Nejdříve stanov, zda je číslo větší, nebo menší než jedna, a teprve potom určuj znaménko exponentu.

\begin{definiceBox}[title={Definice},breakable=false]
\textbf{Vědecký zápis} nenulového reálného čísla je tvar $c\cdot10^k$, kde $k\in\mathbb Z$ a $1\le |c|<10$. Nulu zapisujeme prostě $0$.
\end{definiceBox}


Při posunu desetinné čárky v koeficientu musíme opačně změnit exponent, aby se hodnota nezměnila. Pro kladná čísla stejné jednotky rozhoduje nejprve exponent, při shodě exponentů koeficient. Odhad počítáme před přesným výpočtem; přibližné a přesné rovnosti rozlišujeme symboly $\approx$ a $=$.

\blok{Řešené příklady}

\begin{prikladReseniBox}[title={Příklad 1},breakable=false]
Zapiš vědecky $4\,500\,000$ a $0{,}000072$.
\tcblower
$4{,}5\cdot10^6$ a $7{,}2\cdot10^{-5}$. První koeficient vznikl dělením milionem, druhý násobením stotisícem.
\end{prikladReseniBox}


\begin{prikladReseniBox}[title={Příklad 2},breakable=false]
Vypočítej $\frac{(3\cdot10^5)(4\cdot10^{-3})}{6\cdot10^2}$.
\tcblower
$\frac{12}{6}\cdot10^{5-3-2}=2$. Hrubý odhad $300\,000\cdot0{,}004:600$ také dává hodnotu řádu jednotek.
\end{prikladReseniBox}


\blok{Samostatné procvičení}

\begin{zadaniReseniBox}[title={Procvičení},breakable=false]
\begin{enumerate}[label=\alph*)]
\item Zapiš vědecky $0{,}0036$ a $72\,000$.
\item Uprav $(2\cdot10^4)(3\cdot10^{-6})$.
\item Převeď $2{,}5\cdot10^{-3}\,\mathrm m$ na milimetry.
\end{enumerate}
\tcblower
\begin{enumerate}[label=\alph*)]
\item $3{,}6\cdot10^{-3}$ a $7{,}2\cdot10^4$; koeficient patří do intervalu $\langle1,10)$.
\item $6\cdot10^{-2}=0{,}06$.
\item Protože $1\,\mathrm m=10^3\,\mathrm{mm}$, vyjde $2{,}5\,\mathrm{mm}$.
\end{enumerate}
\end{zadaniReseniBox}


\poznamka{Ověř jednotky i řád výsledku. Řádový odhad zde znamená hrubé určení velikosti, nikoli nevyloženou konvenci „nejbližší mocnina deseti“. Kalkulačku k těmto úlohám nepotřebujeme.}

\blok{Volitelná rezerva}

\begin{prikladReseniBox}[title={Rozšíření procvičení},breakable=false]
Odhadni a potom vypočítej $(7{,}8\cdot10^7):(2{,}6\cdot10^{-2})$.
\tcblower
Odhad $8\cdot10^7:(3\cdot10^{-2})\approx2{,}7\cdot10^9$. Přesně $3\cdot10^9$.
\end{prikladReseniBox}


\end{document}
